% TikZ 图：路端电压变化例题
% Day 20：闭合电路欧姆定律
\begin{tikzpicture}[>=Latex,font=\footnotesize]
  \tikzset{wire/.style={thick}, res/.style={thick}}

  \begin{scope}
    \node[font=\bfseries] at (2.35,2.45) {初态：$R_1=4.5\,\Omega$};
    \draw[wire] (0,0) -- (0,0.45);
    \draw[wire] (0,1.05) -- (0,1.58) -- (0.78,1.58);
    \draw[res] (0.78,1.34) rectangle (1.48,1.82);
    \node[above] at (1.13,1.86) {$r$};
    \draw[wire] (1.48,1.58) -- (2.35,1.58);
    \draw[res] (2.35,1.34) rectangle (3.45,1.82);
    \node[above] at (2.90,1.86) {$R_1$};
    \draw[wire] (3.45,1.58) -- (4.15,1.58) -- (4.15,0) -- (0,0);
    \draw[wire] (-0.24,0.45) -- (0.24,0.45);
    \draw[wire] (-0.36,1.05) -- (0.36,1.05);
    \node[left] at (-0.38,1.05) {$+$};
    \node[right] at (0.34,0.74) {$E$};
    \draw[->,red,thick] (1.70,1.98) -- (2.22,1.98) node[midway,above] {$I_1$};
    \node[blue] at (2.90,0.62) {$U_1=5.4\,\mathrm{V}$};
  \end{scope}

  \begin{scope}[xshift=5.55cm]
    \node[font=\bfseries] at (2.35,2.45) {末态：$R_2=1\,\Omega$};
    \draw[wire] (0,0) -- (0,0.45);
    \draw[wire] (0,1.05) -- (0,1.58) -- (0.78,1.58);
    \draw[res] (0.78,1.34) rectangle (1.48,1.82);
    \node[above] at (1.13,1.86) {$r$};
    \draw[wire] (1.48,1.58) -- (2.35,1.58);
    \draw[res] (2.35,1.34) rectangle (3.45,1.82);
    \node[above] at (2.90,1.86) {$R_2$};
    \draw[wire] (3.45,1.58) -- (4.15,1.58) -- (4.15,0) -- (0,0);
    \draw[wire] (-0.24,0.45) -- (0.24,0.45);
    \draw[wire] (-0.36,1.05) -- (0.36,1.05);
    \node[left] at (-0.38,1.05) {$+$};
    \node[right] at (0.34,0.74) {$E$};
    \draw[->,red,thick] (1.62,1.98) -- (2.30,1.98) node[midway,above] {$I_2$};
    \node[blue] at (2.90,0.62) {$U_2=4\,\mathrm{V}$};
  \end{scope}

  \draw[->,thick] (4.42,1.18) -- (5.24,1.18);
  \node[align=center] at (4.83,0.46) {$R$ 减小\\$I$ 增大\\$U$ 减小};
\end{tikzpicture}
